\scanpage{29}
$n$ assez grand, \unsure{sous} la forme % L1
$M_n(n)$, avec $M_n\in\mathrm{Ob}\,\mathcal M^+(X)$, [pour $n$ fixé, % L2
$M_n$ est bien déterminé par $M$ à isom unique % L3
près, et $\bigl(M_n(n)\simeq M_{n+1}(n+1)\bigr)\Leftrightarrow\bigl(M_{n+1}\simeq M_n(-1)\bigr)$, % L4
on retrouve la description de $\mathcal M(X)$ en % L5
termes de Pseudo $\underset{\longrightarrow}{\mathrm{Lim}}$. ]. Les foncteurs % L6
$T_\ell$ s'étendent à $\mathcal M(X)$, de façon \unsure{ess.|can.} unique, de façon % L7
à rester compatible avec $\otimes$. % L8

7) Dans $\mathcal M(X)$, on a aussi des foncteurs % L9
$\underline{\mathrm{Hom}}$, liés à $\otimes$ par la formule % L10
habituelle % L11
\[
\begin{aligned}
\mathrm{Hom}(P\otimes Q,R)&\simeq\mathrm{Hom}(P,\underline{\mathrm{Hom}}(Q,R))\\ % L12
&\simeq\mathrm{Hom}(Q,\underline{\mathrm{Hom}}(P,R)), % L13
\end{aligned}
\]
et \deleted{\unsure{satisfaisant|satisfaisant par suite}} \inserted{qui s'étendent en} % L14
\[
\mathbb R\mathrm{Hom}(\deleted{\unsure{Q}}P,Q)\qquad P,Q\in\mathrm{Ob}\,D^bM(X) % L15
\]
satisfaisant à la relation analogue % L16
relativement à $\underline{\otimes}$. \uuline{La formation des} % L17
$\underline{\mathrm{Hom}}$ \inserted{et $\mathbb R\underline{\mathrm{Hom}}$} \uuline{est compatible avec les} $T_\ell$. % L18

\deleted{NB.} [NB. On retrouve la formation % L19
des $f^!$ .... ] % L20

% archive sheet number: 28
